\section{Span, Linear Independence, Bases}

We take the following as the working definitions at this point of the course
(Axler / 21-242 style). A \emph{linear combination} of a list
$\cL=(\vv_1,\ldots,\vv_n)$ is a vector $\sum_{i=1}^n a_i\vv_i$ with $a_i\in\F$.

\subsection{Q1. Definition of $\operatorname{Span}(\cL)$}

\begin{promptbox}
State the definition of $\operatorname{Span}(\cL)$.
\end{promptbox}

\begin{coursebox}
\begin{definition*}
If $\cL=(\vv_1,\ldots,\vv_n)$ is a list of vectors in $V$, then
\[
\operatorname{Span}(\cL)
:=\Bigl\{\, \sum_{i=1}^n a_i\vv_i : a_1,\ldots,a_n\in\F \,\Bigr\}.
\]
If $\cL$ is empty, $\operatorname{Span}(\cL)=\{\zero\}$.
\end{definition*}
Equivalently: $\operatorname{Span}(\cL)$ is the set of all linear combinations
of vectors appearing in $\cL$. It is the smallest subspace of $V$ containing
every vector of $\cL$.
\end{coursebox}

\subsection{Q2. Definition of linear independence}

\begin{promptbox}
State the definition of linear independence.
\end{promptbox}

\begin{coursebox}
\begin{definition*}[Axler / sequential form]
A list $(\vv_1,\ldots,\vv_n)$ in $V$ is \emph{linearly independent} if no vector
in the list is a linear combination of the preceding ones: that is,
$\vv_1\neq\zero$ and $\vv_j\notin\operatorname{Span}(\vv_1,\ldots,\vv_{j-1})$
for each $j=2,\ldots,n$. The empty list is linearly independent.
\end{definition*}
A list is \emph{linearly dependent} if it is not linearly independent.

(An equivalent and often more convenient form is Q3 below.)
\end{coursebox}

\subsection{Q3. Independence iff only the trivial combination is $\zero$}

\begin{promptbox}
Prove that $\vv_1,\vv_2,\ldots,\vv_n\subseteq V$ is linearly independent if and
only if $\sum_{i=1}^n a_i\vv_i=\zero$ implies $a_i=0$ for all $i$.
\end{promptbox}

\begin{coursebox}
Write $\cL=(\vv_1,\ldots,\vv_n)$. We use the sequential definition from Q2.

($\Rightarrow$).
Suppose $\cL$ is linearly independent and $\sum_{i=1}^n a_i\vv_i=\zero$.
If some coefficient is nonzero, let $j$ be the \emph{largest} index with
$a_j\neq 0$. Then
\[
\vv_j
=\sum_{i=1}^{j-1}\Bigl(-\frac{a_i}{a_j}\Bigr)\vv_i
\in\operatorname{Span}(\vv_1,\ldots,\vv_{j-1})
\]
(if $j=1$ this says $\vv_1=\zero$). This contradicts sequential independence.
Hence every $a_i=0$.

($\Leftarrow$).
Suppose the only linear combination equal to $\zero$ is the trivial one.
If the list failed sequential independence, some $\vv_j$ would lie in
$\operatorname{Span}(\vv_1,\ldots,\vv_{j-1})$, say
$\vv_j=\sum_{i<j}c_i\vv_i$. Then
\[
\sum_{i<j}c_i\vv_i + (-1)\vv_j+\sum_{i>j}0\cdot\vv_i=\zero
\]
is a nontrivial combination (the coefficient of $\vv_j$ is $-1$), a contradiction.
Thus $\cL$ is linearly independent.
\end{coursebox}

\begin{laterbox}
After bases and dimension: a list of length $n$ in an $n$-dimensional space is
independent iff it is spanning iff it is a basis, and the trivial-combination
criterion is the statement that the map $\F^n\to V$,
$e_i\mapsto\vv_i$, is injective. That is a cleaner packaging, but it uses the
linear-map language from the next section.
\end{laterbox}

\subsection{Q4. Can $\zero$ appear in a linearly independent list?}

\begin{promptbox}
Can $\zero$ appear in a linearly independent list?
\end{promptbox}

\begin{coursebox}
No. If $\vv_j=\zero$, then $1\cdot\vv_j=\zero$ is a nontrivial linear combination
of the list (or, in sequential language, $\zero$ always lies in the span of the
preceding vectors, including the empty span when $j=1$).

So any list containing $\zero$ is linearly dependent.
\end{coursebox}

\subsection{Q5. Every list contains a basis of its span}

\begin{promptbox}
Prove that for any list $\cL\subseteq V$, there exists a sublist $\cL'\subseteq\cL$
such that $\cL'$ is a basis for $\operatorname{Span}(\cL)$.
\end{promptbox}

\begin{coursebox}
Let $\cL=(\vv_1,\ldots,\vv_n)$. Build $\cL'$ by reading $\cL$ from left to right
and keeping a vector if and only if it is not in the span of those already kept.
(Equivalently: drop each vector that is a linear combination of the preceding
vectors of $\cL$.)

\begin{itemize}
\item $\cL'$ is linearly independent by construction (sequential definition).
\item $\operatorname{Span}(\cL')=\operatorname{Span}(\cL)$: every dropped vector
already lay in the span of earlier kept vectors, so by induction each $\vv_j$
lies in $\operatorname{Span}(\cL')$. The reverse inclusion is obvious.
\end{itemize}
A linearly independent spanning list of $\operatorname{Span}(\cL)$ is a basis
of $\operatorname{Span}(\cL)$.
\end{coursebox}

\begin{laterbox}
Once you have the extension lemma and $\dim$: take a maximal linearly independent
sublist $\cL'\subseteq\cL$. Maximality implies it spans $\operatorname{Span}(\cL)$,
else you could add another vector of $\cL$. Finite generation of the span is
what makes maximality work without Zorn.
\end{laterbox}
